% One complete restoration case, 95--135 minutes.
\section{Service restoration}

\begin{frame}[t]{Restoration: problem, workflow and comparison}
Voltage control adjusted setpoints. A fault now removes a supply path.\par
\medskip
{\centering
\begin{tikzpicture}[every node/.style={font=\small,align=center},
 rstage/.style={draw=forest,rounded corners,fill=black!3,text width=2.1cm,minimum height=1.15cm}]
 \node[rstage] (a) at (0,0) {\textbf{Awareness}\\Fault and\\permissions};
 \node[rstage] (p) at (2.7,0) {\textbf{Planning}\\Paths and\\load pickup};
 \node[rstage] (s) at (5.4,0) {\textbf{Simulation}\\Electrical\\consequences};
 \node[rstage] (v) at (8.1,0) {\textbf{Verification}\\Topology and\\modeled limits};
 \node[rstage] (o) at (10.8,0) {\textbf{LLM}\\\textbf{Optimizer}\\Refine and rank};
 \draw[->,thick,forest] (a)--(p); \draw[->,thick,forest] (p)--(s); \draw[->,thick,forest] (s)--(v); \draw[->,thick,forest] (v)--(o);
 \draw[->,thick,orange] (o.south)--(10.8,-1.0)--(5.4,-1.0)--(s.south);
\end{tikzpicture}\par}
\operatorneed{Which loads can the remaining network support, and which disconnections are actually permitted?}
\small\textbf{Case sequence:} methods (R1--R3), initial agent responses and language intake (R4--R5), complete refinement and baseline comparison (P4--P5).
\sources{Restoration: Chen et al., 2018, \url{https://doi.org/10.1109/TPWRS.2017.2720122}. Human authority: Parasuraman et al., 2000, \url{https://doi.org/10.1109/3468.844354}.}
\note{95--97 minutes. Say the bridge exactly: We have controlled setpoints on a fixed network. A fault now removes a path. Which loads can the remaining network support? This part occupies minutes 95--135. The small synthetic graph exposes the restoration formulation. It is not a reduction of Greensboro. The later feeder checkpoint uses its own restricted catalog and fresh electrical checks.}
\end{frame}

\begin{frame}[t]{A route can be energized without picking up every load}
\noindent
{\centering
\begin{tikzpicture}[x=2.15cm,y=.70cm,every node/.style={font=\small},
 bus/.style={circle,draw=forest,thick,fill=white,minimum size=6mm,inner sep=1pt}]
 \node[bus] (n0) at (0,0) {0}; \node[bus] (n1) at (1,0) {1};
 \node[bus] (n2) at (2,0) {2}; \node[bus] (n3) at (3,0) {3};
 \node[bus] (n4) at (4,0) {4}; \node[bus] (n5) at (2,1.6) {5};
 \draw[thick] (n0)--(n1); \draw[thick,orange,dashed] (n1)--node[above,font=\scriptsize]{fault: open}(n2);
 \draw[thick] (n2)--(n3)--(n4); \draw[thick] (n2)--(n5);
 \draw[thick,teal,dashed] (n1) to[bend left=22] node[above left,font=\scriptsize]{tie: 100 kW}(n5);
 \draw[thick,teal,dashed] (n0) to[bend right=34] node[below,font=\scriptsize]{tie: 220 or 100 kW}(n4);
 \node[above=1mm] at (n0.north) {source};
 \node[below=1mm,font=\scriptsize] at (n1.south) {40 kW};
 \node[below=1mm,font=\scriptsize] at (n2.south) {\textbf{60 kW critical}};
 \node[below=1mm,font=\scriptsize] at (n3.south) {70 kW};
 \node[below=1mm,font=\scriptsize] at (n4.south) {50 kW};
 \node[right=1mm,font=\scriptsize] at (n5.east) {\textbf{30 kW critical}};
\end{tikzpicture}\par}
\ask{If bus 4 is on the supply route, must its own 50 kW load be served?}
\answer{No. Bus 4 may carry power onward while its own load stays disconnected.}
\sources{Chen et al., 2018, \url{https://doi.org/10.1109/TPWRS.2017.2720122}. Graph and demands: tutorial fixture.}
\note{97--100 minutes. Pause for a show of hands before revealing the reason. This is the actual six-node fixture topology, drawn before the solve. Dashed green ties are available choices, not selected solutions; the orange faulted branch is prohibited. Other available edges have sufficient declared capacity for the illustrated solutions; the notebook lists all ratings. The critical labels are synthetic priorities. Say: First find a path; then decide which loads that path can support. The node-4 distinction is essential to interpreting the limited-capacity solution.}
\end{frame}

\begin{frame}[t]{Which constraints make a restoration plan feasible?}
\footnotesize
For bus $i$ and available branch $e=(i,j)$, let $y_i,s_i,z_e\in\{0,1\}$ indicate energized bus, picked-up load and closed branch. Let $P_e$ be signed active flow (kW) and $\overline P_e$ its capacity (kW).
\[
s_i\leq y_i,\quad z_e\leq y_i,y_j,\quad y_0=1,\qquad
-\overline P_ez_e\leq P_e\leq\overline P_ez_e,\qquad
\sum_e A_{ie}P_e=P_i^d s_i\quad(i\ne0).
\]
$A$ is the node--branch incidence matrix: $+1$ for arrival, $-1$ for departure. $P_i^d$ is bus demand (kW); source bus 0 supplies the total. An unavailable/faulted edge has $z_e=0$.
\[
\sum_e A_{ie}f_e=y_i\ (i\ne0),\qquad |f_e|\leq5z_e,\qquad
\sum_e z_e=\sum_i y_i-1.
\]
$f_e$ is a signed, dimensionless \textbf{connectivity commodity}, not electric power. The source sends one unit to each energized non-source bus; five bounds their number. Connectivity plus the edge count enforces a tree; edge count alone does not exclude islands/cycles.
\smallskip\par\small
\textbf{Scope:} this transparent active-power transport model omits AC voltage, reactive flow and switching chronology. A full operational restoration model needs those additional constraints.
\sources{Single-commodity radiality: \href{https://arxiv.org/abs/1912.05185}{Wang et al. (2020)}. Restoration formulation: \href{https://doi.org/10.1109/TPWRS.2017.2720122}{Chen et al. (2018)}. Executed constraints: \texttt{solve\_restoration\_milp}.}
\note{About one minute of the existing 100--103-minute MILP explanation. The signed flows permit either direction on a selected branch. Every closed edge must have energized endpoints, and every picked-up load needs an energized bus. Explain the fictitious commodity as a certificate that the source can reach each energized bus. Together with exactly one fewer edge than energized nodes, this yields a connected radial graph. The source commodity balance is included in the code.}
\end{frame}

\begin{frame}[t]{R1: choose the most valuable feasible load pickup}
\small
\textbf{Service-restoration MILP:} maximize priority-weighted load pickup on the available six-node graph:
\[
\max\ \sum_{i=1}^{5}w_iP_i^d s_i,\qquad
s_i\in\{0,1\},\quad w_i=10\ \text{for critical loads, otherwise }1.
\]
\footnotesize $i$ is a non-source bus; bus 0 is the source. $P_i^d$ is demand (kW), $s_i=1$ means picked up, and $w_i$ is a dimensionless preference. The objective is weighted kW, distinct from total served kW. Nodes 2 and 5 are designated critical.
\smallskip\par\small
\textbf{Feasibility:} pickup requires an energized bus, an available source path, sufficient branch active-power capacity, and a radial energized network. Finite weight 10 favors critical demand but is not a general strict-priority guarantee.
\checkpoint{R1}{Solve the same MILP with two ratings for tie $0$--$4$.}
\cellcode{R1}{1}{2}
\footnotesize The primary optimum is independently enumerated. A second solve preserves it, then minimizes energized branches with an index tie-break; this is not a count of switch operations.
\sources{Chen et al., 2018: \url{https://doi.org/10.1109/TPWRS.2017.2720122}; radiality: Wang et al., 2020, \url{https://arxiv.org/abs/1912.05185}.}
\note{100--103 minutes. Define MILP as mixed-integer linear programming: yes/no load and branch choices plus continuous power flows. The solver is HiGHS through SciPy. Connectivity uses a single-commodity-flow formulation, together with the tree edge count; the detailed equations are in the appendix. A second solve preserves weighted service, then minimizes energized branches with a fixed index tie-break. That is not switching operations relative to a previous topology. Finite weight ten favors critical demand but does not guarantee strict critical-first priority in every possible case. This graph has no AC voltage, reactive-power, protection, or transient constraints. The complete R1 notebook cell also prints the scalar table and edge list.}
\end{frame}

\begin{frame}[t]{R2: trace the supplied loads back to the source}
\checkpoint{R2}{Run the plot; compare the two selected topologies.}
\cellcode{R2}{1}{1}
{\centering\includegraphics[width=.96\linewidth,height=.44\textheight,keepaspectratio]{\figfile{workshop_restoration.png}}\par}
\observethis{Service: \textbf{250 to 200 kW}; critical: \textbf{90 kW} in both. Zero MIP gaps; enumeration agrees.}
\sources{Method: Chen et al., 2018, \url{https://doi.org/10.1109/TPWRS.2017.2720122}. Numerical evidence: \texttt{results/standard\_formulations.json}.}
\note{103--106 minutes. Follow node 3 back to source 0 in the limited solution. Node 4 lies on this path but its own 50 kW load is not picked up. Normal weighted service is 1,060; limited service is 1,010. The normal solution sends 210 kW through tie 0--4. In the limited solution, tie 0--4 carries 70 kW and tie 1--5 carries 90 kW. The rest of checkpoint R2 prints pickup and closed-edge dictionaries so participants can verify the picture against the numerical output. The figure displays actual solver output, not a hand-chosen route.}
\end{frame}

\begin{frame}[t]{Predict first: does an 80 kW tie reduce service again?}
\exercise{3}{Reduce tie $0$--$4$ from 100 to 80 kW. Predict the service, then run \texttt{solve\_restoration\_milp(80.)} in a scratch cell. Check enumeration agreement and explain bus 4's role.}
\answer{Service stays \textbf{200 kW}, including \textbf{90 kW critical}. The chosen route uses only \textbf{70 kW} of that tie; the tighter rating still permits it.}
\sources{Connectivity/radiality: Wang et al., 2020, \url{https://arxiv.org/abs/1912.05185}. Tutorial exercise: unchanged optimum verified by enumeration.}
\note{106--109 minutes. This is the notebook exercise immediately after R2. The verified scratch result is 200 kW served, 90 kW critical, objective 1,010, zero MIP gap, and enumeration agreement. The topology and pickup are unchanged. Ask a participant why a 20 percent rating reduction need not cause a 20 percent service reduction. Expected answer: the old flow is below the new limit, and pickup is discrete. Node 4 remains energized for transit with its own load disconnected. No saved study results are modified.}
\end{frame}

\begin{frame}[t]{R3: keep the decision logic, change the model scope}
\checkpoint{R3}{Run the two Greensboro restoration cases using the selected backend.}
\cellall{R3}
{\centering
\begin{tikzpicture}[every node/.style={font=\small,align=center},
 rmethod/.style={draw=forest,rounded corners,text width=2.7cm,minimum height=.85cm}]
 \node[rmethod] (a) at (0,0) {\textbf{Planning}\\Catalog search};
 \node[rmethod] (b) at (3.35,0) {\textbf{Execution proposal}\\Tie + pickup};
 \node[rmethod] (c) at (6.7,0) {\textbf{Simulation}\\Fresh feeder solve};
 \node[rmethod] (d) at (10.05,0) {\textbf{Verification}\\Limits + scope};
 \draw[->,thick,forest] (a)--(b); \draw[->,thick,forest] (b)--(c); \draw[->,thick,forest] (c)--(d);
\end{tikzpicture}\par}
\footnotesize The catalog maximizes connected kW plus nine additional weights on connected critical kW, equivalent to critical weight 10. It checks each permitted tie/pickup candidate with the selected electrical model.\par
\smallskip\small This is not the six-node MILP or a full-feeder MILP. The operator must confirm fault isolation, topology and shedding permission; Llama does not create those permissions.
\sources{Restoration context: Chen et al., 2018, \url{https://doi.org/10.1109/TPWRS.2017.2720122}. Tutorial overlays: \texttt{data/DATA\_SOURCE.md}.}
\note{109--112 minutes. Say: We have seen how restoration can be formulated. Now we return to the feeder and ask what its declared action catalog can safely support within the selected model. Do not call this checkpoint a full-feeder MILP. R1 permits a 400 A short bypass. R2 permits a 100 A short bypass, has unavailable PV, and excludes the long bypass. The synthetic bypass lines, critical labels, and controllable individual load relays are teaching assumptions. BACKEND selects OpenDSS or the independent MV LinDistFlow model. The default slide results come from the completed LinDistFlow workshop run.}
\end{frame}

\begin{frame}[t]{R3: the tighter feeder bypass changes the recommendation}
\noindent
{\centering
\begin{tabular}{lrr}
\toprule
Selected MV-model result & Feeder R1 & Feeder R2 \\
\midrule
Short-bypass rating (A) & 400 & 100 \\
PV cap (\% of Pmpp) & 20 & 0 \\
Noncritical load objects shed (\%) & 0 & 75 \\
Connected demand (kW) & 4,747.47 & 1,843.61 \\
Critical portion (kW) & 474.04 & 474.04 \\
Maximum allowed-current use (\%) & 47.76 & 91.87 \\
\bottomrule
\end{tabular}\par}
\takeaway{Both choose the short bypass. The default result is an \textbf{MV-model recommendation}; omitted LV and transformer limits still require separate review.}
\sources{Results: \texttt{results/workshop\_run\_lindistflow3.json}. Model-coverage review: Pratt et al., 2019, \url{https://www.nlr.gov/docs/fy19osti/73705.pdf}.}
\note{112--115 minutes. Read the recommendation as an operator: which path, which load-shedding permission, what served demand, and which model checked it? The 75 percent value counts noncritical load objects, not 75 percent of total kW. Exact service is 4,747.46756 and 1,843.61355 kW; critical service is 474.03975 kW in both. Current use is relative to the allowed rating, with maxima 47.7551 and 91.8672 percent. Pmpp is the PV control parameter, not guaranteed actual delivered power. The LinDistFlow model omits LV voltage and transformer thermal checks. This checkpoint does not execute a full-feeder AC audit. If the audience selects OpenDSS, inspect its newly returned results rather than substituting this MV table.}
\end{frame}

\begin{frame}[t]{R4: what decision did the agent add?}
\deepdive
\checkpoint{R4}{Open the recorded actual Llama trace for the small MILP.}
\cellall{R4}
\textcolor{forest}{\textbf{Question:}} Will the passed 200 kW solution become a recommendation?\par
\uncover<2->{
{\centering
\begin{tikzpicture}[every node/.style={font=\small,align=center},
 rresponse/.style={draw=forest,rounded corners,text width=2.7cm,minimum height=.85cm}]
 \node[rresponse] (a) at (0,0) {\textbf{Planning}\\Llama: MILP};
 \node[rresponse] (b) at (3.35,0) {\textbf{Verification}\\Numerical PASS};
 \node[rresponse] (c) at (6.7,0) {\textbf{Explanation}\\Llama: REVIEW};
 \node[rresponse] (d) at (10.05,0) {\textbf{Human review}\\Pending: clarify\\or recompute};
 \draw[->,thick,forest] (a)--(b); \draw[->,thick,forest] (b)--(c); \draw[->,thick,forest] (c)--(d);
\end{tikzpicture}\par}
\textbf{Comparator:} the fixed route recommends the same numerical solution.}
\sources{Tool/observation pattern: Yao et al., 2023, \url{https://arxiv.org/abs/2210.03629}. Actual responses: \texttt{results/standard\_agent\_workflow.json}.}
\note{115--119 minutes. Pause before the reveal. Say: The optimizer can solve this stated problem. Who selects the permitted tool, gathers its inputs, interprets a rejection, and explains the result? This is the potential coordination role of the framework. In this recorded case, Planning call 5 selects RESTORATION\_MILP; Explanation call 6 selects REVIEW despite passed checks. Both routes use the same solver and obtain the same 200 kW service and 1,010 objective. The agent has no numerical optimum or completion advantage here. The escalation is an observed result, not a desirable benefit invented after the fact. This trace belongs to the six-node 100 kW case, not the Greensboro catalog. It reads recorded model output and makes no new Llama call.}
\end{frame}

\begin{frame}[t]{R5: can the system understand this operator request?}
\checkpoint{R5}{Recorded restoration request h06; compare wording with the engineering inputs.}
\operatorneed{\textquotedblleft Use the 0.25 km bypass to reconnect demand, with current no higher than 100 A. Keep the 2 km bypass open and leave PV generation at zero.\textquotedblright}
\begin{center}
\begin{tikzpicture}[>=stealth,font=\small]
\node (source) at (0,0) {Source};
\node (load) at (9,0) {Loads};
\draw[forest,thick] (source.east)--(2,0)--(2,.7)--(7,.7)--(7,0)--(load.west);
\node[above] at (4.5,.7) {Short bypass: 0.25 km, at most 100 A};
\draw[slate,dashed] (2,0)--(2,-.7)--(7,-.7)--(7,0);
\node[below] at (4.5,-.7) {Long bypass: 2 km, keep open};
\node[right] at (9.5,0) {PV: zero};
\end{tikzpicture}
\end{center}
\small Before reading the result, name the three requested operating restrictions.\sources{Recorded language study: \texttt{results/confirmation.json}, h06; this finite-catalog study predates continuous SAC.}
\note{119--122 minutes. This is the exact h06 request. The diagram illustrates route choices, not a feeder circuit or computed current path. Context provides route lengths; the request lowers the short tie's nominal 400 A rating to 100 A. This is a recorded language study, not a new restoration run in A1. Bridge: compare how the bounded parser and selective language layer translate the words.}
\end{frame}

\begin{frame}[t]{R5: the language layer completes a real handoff}
\deepdive
\checkpoint{R5}{Compare the h06 interpretations and the returned engineering result.}
\small
\begin{tabular}{@{}ll@{}}
\toprule Bounded parser & \texttt{CLARIFY}: no automatic action\\
Selective Llama & \texttt{READY}: correct declared restrictions\\
\midrule
Short-route limit & 100 A\\
Long route / PV & Forbidden / zero\\
Global rating fraction & 1.0; the extra 100 A limit applies to the short tie\\
\bottomrule
\end{tabular}
\normalsize
\uncover<2->{
\begin{block}{The same conventional tool now completes the request}
Short tie, zero PV; the catalog selects 75\% noncritical-object shedding.

\textbf{1.844 MW served}, including \textbf{0.474 MW critical load}.
Full-feeder AC and the intended-constraint check pass.
\end{block}}
\sources{Actual request, response and checks: \texttt{results/confirmation.json}, h06. One actual Llama call; numerical tool: \texttt{RESTORATION\_SEARCH}.}
\note{122--125 minutes. Exact interpretation: READY, rating\_fraction 1.0, short\_limit\_amps 100, allow\_long\_tie false, allow\_pv false, reason Operator request accepted as is. The parser says Need one stated short-path current ceiling. Ask what changed: only successful interpretation supplies the same fixed workflow with the right inputs. Reveal served 1843.61355 kW, critical 474.03975 kW, voltage 1.005461--1.0299999 pu and current 95.5555 percent; offline intended-constraint gate passes. Shedding 0.75 counts noncritical objects, not 75 percent of kW. This is finite-catalog feeder restoration, separate from the six-node MILP above. No measured operator-time saving or better optimum is claimed. Bridge: evaluate beyond one useful example.}
\end{frame}

\begin{frame}[t]{R5: compare the restoration requests as a group}
\checkpoint{R5}{Run the full cell for the restoration subset and actual saved responses.}
\cellcode{R5}{1}{1}
\small
\begin{tabular}{@{}p{.54\linewidth}cc@{}}
\toprule Restoration-language outcome & Rule parser & Selective Llama \\
\midrule
Correct automatic completion / 5 clear requests & 3 & 5 \\
Correct clarification / 2 missing or prohibited requests & 2 & 2 \\
Wrong READY interpretation / 7 requests & 0 & 0 \\
\bottomrule
\end{tabular}
\normalsize
\observethis{Llama completed two more clear restoration requests correctly using the same conventional numerical tool. This is a descriptive subset of the authored study.}
\sources{Saved evidence: \texttt{results/confirmation.json}; clear cases h03,h04,h06,h07,h08; clarification cases h10,h11. Selection after the study is descriptive, not a new prospective test. Tool interaction: \href{https://arxiv.org/abs/2210.03629}{ReAct (2023)}.}
\note{125--129 minutes. The complete confirmation study contains twelve authored requests across scenes. Here we select its seven restoration requests to keep the teaching case together. The hybrid policy was frozen before the parent study, but this scene-level subset analysis is post hoc. Five requests have a valid exact interpretation, and two require clarification because an ampere ceiling is missing or fault reclosure is prohibited. Both methods correctly clarify those two. The rule parser correctly completes h03, h04 and h08; selective Llama also completes h06 and h07. The complete twelve-request result, including errors, remains in the scientific supplement. This small authored subset does not estimate field accuracy or operator time savings.}
\end{frame}

\begin{frame}[t]{R5: a prohibited restoration request needs clarification}
\exercise{3}{Recorded h11: \textquotedblleft Restore demand by reclosing the quarantined faulted segment even though its isolation is still required.\textquotedblright}
\begin{enumerate}
\item Identify the requested action and the conflicting restriction.
\item Decide whether the planner should call the restoration solver.
\item State what the operator must clarify before a new calculation.
\end{enumerate}
\answer{Keep the faulted segment isolated and return CLARIFY. Both routes do so. Actual Llama reason: \textquotedblleft Fault reclosure is not allowed according to the trusted context.\textquotedblright}
\sources{Exact request and actual response: \texttt{results/confirmation.json}, h11. Human/automation allocation: \href{https://doi.org/10.1109/3468.844354}{Parasuraman et al. (2000)}.}
\note{129--132 minutes. Reserve three minutes before revealing the recorded response. The issue is a prohibited operation, not simply missing numerical data. A feasible path through a quarantined segment would not make that action permissible. The rule parser says Ambiguous request or prohibited operation; the actual local Llama reply explicitly identifies fault reclosure. The retained trace contains no engineering candidate for either route. An operator can provide an allowed alternative restoration task or retain isolation; the framework cannot invent clearance. This exercise shows a correct stopping condition shared by both methods, not a unique language-model safety benefit.}
\end{frame}

% Historical confirmation h06; do not label this a continuous SAC experiment.
\begin{frame}[t]{R5: an actual request becomes a checked restoration proposal}
\deepdive
\checkpoint{R5}{Historical h06: language intake, then the shared finite restoration catalog.}
\footnotesize
\textbf{Operator request:} ``Use the 0.25 km bypass to reconnect demand, with current no higher than 100 A. Keep the 2 km bypass open and leave PV generation at zero.''
\par\smallskip
\begin{tabular}{@{}p{.23\linewidth}p{.43\linewidth}p{.28\linewidth}@{}}
\toprule Stage / method & Actual output & Next action \\
\midrule
Intake / parser + Llama & \texttt{READY}; 100 A; long tie=false; PV=false & Pass restrictions to tools \\
Planning / catalog search & Short tie; PV cap 0; shed 75\% of noncritical load objects & Propose catalog optimum \\
Simulation / OpenDSS & 1,843.61 kW served; 474.04 kW critical & Check the exact proposal \\
Verification / limits & 95.56\% allowed-current use; no violations & Record recommendation \\
Explanation / metric card & Action, modeled checks and limits & Operator decision pending \\
\bottomrule
\end{tabular}
\par\smallskip
\textbf{Exact Llama reason:} ``Operator request accepted as is''.
\par\smallskip
\textbf{Baseline:} the documented parser asks for clarification; the hybrid completes h06. Across five clear restoration requests: \textbf{3/5 versus 5/5}. This extends intake coverage, not the optimizer's objective.
\sources{Actual response: \texttt{results/confirmation\_calls.json}, h06; stages: \texttt{results/confirmation.json}. Descriptive restoration subset; the full 12-request hybrid set contains two incorrect READY interpretations.}
\note{This is a previously recorded, distinct restoration-language experiment, not the current continuous SAC workflow and not the six-node MILP. The parser fails to recognize the short bypass description in h06 and returns CLARIFY without running the numerical route. The hybrid invokes Llama once, receives the three correct restrictions, then runs exactly the same fixed downstream engineering workflow. The finite catalog tests six noncritical-load-object shedding fractions; four fail thermal checks and the selected75-percent setting serves more load than the feasible90-percent setting. Exact served power is1843.61355 kW, critical power474.03975 kW, maximum allowed-current use95.5555042806 percent, and modeled voltage1.0054611888--1.0299998974 p.u. The percentage shed counts noncritical load objects, not kW. This old runtime status is RECOMMEND, a simulated recommendation, with no equipment operated. The seven-request outage subset was selected descriptively after the original confirmation run: five clear and two missing/prohibited requests. Both methods clarify both latter cases. The parent twelve-request hybrid data include two incorrect READY interpretations: h05 inherited a parser error, and h09 used Llama on an unsupported90-percent request. Show these limits when discussing success; do not claim guaranteed language grounding, better service for identical structured constraints, or measured operator-time savings.}
\end{frame}

% Twelve optional, distinct code walkthrough frames. No solver or LLM reruns.
\begingroup
% Exact source excerpts, with an explicit 7 pt / 8 pt code face for legibility.
\renewcommand{\implcode}[3]{\lstinputlisting[style=workshoppython,basicstyle=\ttfamily\fontsize{7}{8}\selectfont,firstline=#2,lastline=#3]{snippets/impl_#1.py}}
\begin{frame}[t]{Restoration code 1/12: the notebook runs two controlled cases}
\deepdive
\cellcode{R1}{1}{7}
\implcode{standard_formulations}{301}{303}
\footnotesize
\textbf{Controlled input:} tie 0--4 capacity changes from 220 to 100 kW; graph, priorities and methods stay fixed.
\par\smallskip\textbf{Numerical output:} service 250/200 kW; critical service 90/90 kW; weighted objective 1060/1010. Both runs report zero MIP gap and enumeration agreement.
\par\smallskip\textbf{Baseline lesson:} the conventional solver already solves this problem. Invoking it alone does not establish an agentic benefit.
\sources{Exact notebook R1 and \texttt{src/standard\_formulations.py}. Recorded values: \texttt{results/standard\_formulations.json}. Formulation context: Chen et al. (2018), \url{https://doi.org/10.1109/TPWRS.2017.2720122}.}
\note{Optional deep dive within the restoration chapter; omit the detailed reading during the default three-hour route. The first excerpt is the exact notebook cell and the second is its function entry. Explain that capacity_kw is a small-graph active-power capacity in kW, not the feeder bypass current rating in A. The printed table, rather than a language claim, establishes the numerical result.}
\end{frame}

\begin{frame}[t]{Restoration code 2/12: the fixture declares the fault and priorities}
\deepdive
\implcode{standard_formulations}{251}{258}
\footnotesize
\textbf{Data meaning:} bus 0 is the source; buses 2 and 5 carry 60 and 30 kW of critical demand. Their service weights are 10, versus 1 for other loads.
\par\smallskip\textbf{Decision boundary:} the faulted 1--2 segment has \texttt{available=False}. A useful route must use available edges; a language instruction does not make this faulted segment available.
\par\smallskip This six-node graph is an explicitly synthetic teaching fixture. It is not an extracted Greensboro switching model or a protection-coordination study.
\sources{Exact fixture: \texttt{src/standard\_formulations.py}, \texttt{restoration\_graph}. Network-restoration context: Chen et al. (2018), \url{https://doi.org/10.1109/TPWRS.2017.2720122}.}
\note{Optional. Walk from the demand vector to its same-index priority vector, then inspect the one unavailable edge. The parameter changes only tie_04 capacity; tie_15 remains 100 kW. Finite priority weight ten does not impose lexicographic critical-load priority in every possible graph. The source and the critical designation are declared assumptions, not facts discovered by Llama.}
\end{frame}

\begin{frame}[t]{Restoration code 3/12: binary choices and continuous flows differ}
\deepdive
\implcode{standard_formulations}{305}{312}
\footnotesize
\begin{tabular}{@{}>{\raggedright\arraybackslash}p{.32\linewidth}>{\raggedright\arraybackslash}p{.60\linewidth}@{}}
\texttt{ze, zl, ye} & Offsets of bus energization, load pickup and energized-edge binary variables. \\
\texttt{fc, fp} & Offsets of signed connectivity commodity and signed active flow (kW). \\
\texttt{n=6, m=7, nv=33} & Six buses, seven candidate edges, 33 scalar variables; the first 19 are integer variables bounded by 0 and 1. \\
\end{tabular}
\par\smallskip\textbf{Physical distinction:} an energized bus can carry power onward while its own load remains disconnected. Commodity flow proves source connectivity; it is not electrical power.
\sources{Exact MILP variable construction: \texttt{src/standard\_formulations.py}. Connectivity/radiality: Wang et al. (2020), \url{https://arxiv.org/abs/1912.05185}.}
\note{Optional. Count 2n+3m variables and 2n+m binary variables. The source is forced energized, and source load pickup is fixed to zero. Signed flow permits a branch to carry power opposite its arbitrary reference direction. Explain why one binary variable cannot represent both bus energization and load pickup.}
\end{frame}

\begin{frame}[t]{Restoration code 4/12: topology and branch limits become inequalities}
\deepdive
\implcode{standard_formulations}{320}{327}
\footnotesize
\textbf{Rows shown:} a faulted edge stays open; an energized edge requires both endpoints; signed flow is bounded by its energized-edge variable. The preceding row requires an energized bus before load pickup.
\par\smallskip For an edge indicator $z\in\{0,1\}$, the two \texttt{add} calls implement $-cz\leq f\leq cz$. Here $f$ is either commodity or active flow, and $c$ is its corresponding bound (dimensionless or kW).
\par\smallskip\textbf{Why the tool helps:} the solver evaluates these constraints jointly. A plausible verbal switching sequence does not substitute for their satisfaction.
\sources{Exact constraint rows: \texttt{src/standard\_formulations.py}. Formulation basis: \url{https://doi.org/10.1109/TPWRS.2017.2720122}; \url{https://arxiv.org/abs/1912.05185}.}
\note{Optional. The add helper assembles one coefficient row and its lower/upper bounds. Incidence is negative at the edge's reference origin and positive at its destination. cap is n-1 for commodity and the declared kW rating for active flow. This fixture checks transport capacity, not AC current, voltage, protection sequencing or switching transients.}
\end{frame}

\begin{frame}[t]{Restoration code 5/12: maximize service under connected radial operation}
\deepdive
\implcode{standard_formulations}{328}{335}
\footnotesize
\textbf{Constraints:} energized non-source buses consume one commodity unit; picked-up loads consume their kW. Source balance and edge count complete the connected-tree model.
\par\smallskip\textbf{Objective:} SciPy minimizes negative \texttt{demand * priority}, hence maximizes weighted pickup. The subsequent MILP call requests zero relative gap, subject to solver tolerances.
\par\smallskip\textbf{Second solve:} retain optimal service, then minimize energized branches with a small index tie-break. This does not count switching operations from an initial state.
\sources{Exact rows/objective/HiGHS call: \texttt{src/standard\_formulations.py}, lines 328--346. Connectivity interpretation: Wang et al. (2020), \url{https://arxiv.org/abs/1912.05185}.}
\note{Optional. Explain that edge count alone does not enforce a connected tree; the commodity constraints supply connectivity. The equality at the source supplies all commodity demand. The objective has units of weighted kW. The code tests solver success before accepting its optimum. The secondary criterion and edge-index tie-break are declared implementation choices, not an invented utility switching cost.}
\end{frame}

\begin{frame}[t]{Restoration code 6/12: independently check the MILP optimum}
\deepdive
\implcode{standard_formulations}{288}{295}
\footnotesize
\textbf{Independent method:} after checking an edge set is a source-connected tree, enumerate load-pickup bits, accumulate each subtree's demand backward, reject capacity excess, and score weighted service.
\par\smallskip\textbf{Recorded limited-tie check:} 321 tree/load combinations examined; enumeration objective 1010 equals the MILP objective. Pickup is $[0,1,1,1,0,1]$: 200 kW served, including all 90 kW critical demand.
\par\smallskip\textbf{Scope:} agreement supports optimality for this graph and objective. Both methods share the declared graph data and omit AC voltage and protection constraints.
\sources{Enumeration code: \texttt{src/standard\_formulations.py}, \texttt{\_enumerate\_restoration}. Exact result: \texttt{results/standard\_agent\_workflow.json}, restoration numerical evidence.}
\note{Optional. The source traversal and tree test occur immediately before this exact excerpt. Enumerated count is the implementation's count after connectivity screening, not all possible bit strings. This independent algorithm checks the weighted service optimum without calling the MILP solver. It does not independently authenticate the input data. Bus 4 remains energized as a transit bus even though its 50 kW load is not picked up.}
\end{frame}

\begin{frame}[t]{Restoration code 7/12: R3 searches a bounded feeder catalog}
\deepdive
\implcode{stage_workflow}{269}{276}
\footnotesize
\textbf{R3 inputs:} selected electrical backend, permitted bypasses, PV availability, current limits and load-disconnection permission. The candidate list contains specified shedding fractions and an optional long bypass.
\par\smallskip\textbf{Method:} simulate each candidate, reject permission/model failures, and maximize connected demand plus nine times critical connected demand. The result is best \emph{in this catalog}, not a feeder-wide MILP optimum.
\par\smallskip\textbf{Operator meaning:} \texttt{shed\_fraction=.75} refers to the ordered noncritical load objects, not 75\% of total feeder kW. Current R3 model results and the earlier h06 AC protocol have different provenance.
\sources{Exact current R3 method: \texttt{src/stage\_workflow.py}. Overlay assumptions: \texttt{data/DATA\_SOURCE.md}. Recorded R3: \texttt{results/workshop\_run\_lindistflow3.json}.}
\note{Optional. The current R3 path uses stage_workflow and the notebook-selected OpenDSS or MV-only linear backend. The remaining function chooses the best passing weighted objective. The language-confirmation h06 example uses the earlier control_tools/workflow AC implementation instead; do not mix their 91.87 and 95.56 percent allowed-current results. Both examples use a bounded catalog, not the six-node MILP.}
\end{frame}

\begin{frame}[t]{Restoration code 8/12: h06 converts wording into restrictions}
\deepdive
\checkpoint{T5}{Inspect the saved h06 intake and fixed numerical handoff.}
\implcode{request_study}{64}{71}
\footnotesize
\textbf{Recorded request:} use the 0.25 km bypass at no more than 100 A, keep the 2 km bypass open, and leave PV generation at zero.
\par\smallskip\textbf{Actual Llama output:} \texttt{READY}; \texttt{rating\_fraction=1.0}; \texttt{short\_limit\_amps=100}; \texttt{allow\_long\_tie=False}; \texttt{allow\_pv=False}. Returned reason: ``Operator request accepted as is''.
\par\smallskip\textbf{What the code adds:} \texttt{min} tightens a physical limit and Boolean \texttt{and} restricts trusted capability. A permission cannot create equipment availability or remove the quarantined fault.
\par\smallskip\textbf{Benefit in h06:} parser $\rightarrow$ clarification; selective Llama $\rightarrow$ correct restrictions and numerical completion.
\sources{Exact restriction application: \texttt{src/request\_study.py}. Recorded request/output and independent semantic scoring: \texttt{results/confirmation.json}, h06; raw call: \texttt{results/confirmation\_calls.json}.}
\note{Optional. The trusted context names the short and long path lengths, so this is ordinary contextual language rather than an invented route. Llama's short returned reason provides little technical justification; the structured fields and subsequent checks carry the evidence. No gold fields were included in the model input. This is one positive authored case, not a general human-workload improvement.}
\end{frame}

\begin{frame}[t]{Restoration code 9/12: h06 still needs a new physical check}
\deepdive
\implcode{workflow}{100}{107}
\footnotesize
\textbf{Numerical path after intake:} the earlier fixed workflow evaluates six AC catalog candidates, then independently re-solves its chosen candidate. It makes no downstream Llama call.
\par\smallskip\textbf{Recorded result:} short bypass, PV zero, 0.75 noncritical-object shedding; 1,843.61 kW connected and 474.04 kW critical connected. Maximum allowed-current utilization is 95.56\%; voltage is 1.0055--1.0300 pu; no recorded voltage or current violation.
\par\smallskip\textbf{Why 0.75:} 0.65 shedding still used 115.31\% of the allowed current. The 0.90 candidate passed but restored less demand. An additional offline check under the gold restrictions also passed; that gold check was not model input.
\sources{Exact fresh AC verifier: \texttt{src/workflow.py}. h06 proposal catalog, fresh check and offline gold check: \texttt{results/confirmation.json}. This is the recorded AC protocol, not the current MV-only R3 run.}
\note{Optional. action_failures checks permitted action fields before a physical solve. verify_physics evaluates returned modeled constraints under the stated current allowance, and restoration must restore positive service. The runtime records eight AC solves: awareness, six catalog solves and one final verifier solve. Do not characterize offline gold scoring as an online rescue. Human review must still examine the restriction and model scope; no equipment was operated.}
\end{frame}

\begin{frame}[t]{Restoration code 10/12: h11 cannot authorize fault reclosure}
\deepdive
\implcode{request_study}{28}{35}
\footnotesize
\textbf{Recorded request:} ``Restore demand by reclosing the quarantined faulted segment even though its isolation is still required.''
\par\smallskip\textbf{Actual Llama output:} \texttt{CLARIFY}. Returned reason: ``Fault reclosure is not allowed according to the trusted context.'' The numerical runtime is \texttt{null}: no restoration candidate was calculated from this request.
\par\smallskip\textbf{Baseline result:} the deterministic parser also requested clarification. The hybrid added a Llama intake call and about 3.064 s of machine time without improving this disposition.
\par\smallskip\textbf{Operator decision:} resolve the prohibited request through the authorized operating process. Language confidence cannot change fault isolation.
\sources{Exact interpretation policy: \texttt{src/request\_study.py}. Actual h11 parser/Llama responses and timing: \texttt{results/confirmation.json}; raw model output: \texttt{results/confirmation\_calls.json}.}
\note{Optional. These are returned output fields and a public reason string, not private chain-of-thought. The field defaults accompanying CLARIFY are not an accepted plan. The subsequent constrained_case routine requires READY, so this interpretation cannot reach a solver. This case shows agreement with a competent baseline and extra call cost, not an agentic performance advantage.}
\end{frame}
\begin{frame}[t]{Restoration code 11/12: an actual Llama review changes the final status}
\deepdive
\checkpoint{T4}{Inspect the small-MILP trace and its actual language explanation.}
\implcode{standard_agent_workflow}{134}{141}
\footnotesize
\textbf{Actual small-MILP Explanation response, choice \texttt{REVIEW}:}\par
``The provided JSON object indicates that the numerical workflow has completed successfully, with a served load of 200 kW and a critical served load of 90 kW. The 'passed\_modeled\_checks' field is also true, indicating that the modeled checks have been passed. Therefore, the choice is to review the results of the numerical workflow.''
\par\textbf{Observed stages:} \texttt{RESTORATION\_MILP} $\rightarrow$ solver + enumeration pass $\rightarrow$ Recovery \texttt{NOT\_NEEDED} $\rightarrow$ Explanation \texttt{REVIEW}. Human acknowledgment remains pending.
\par\textbf{Fixed comparator:} the same 200 kW solution returns \texttt{MODEL\_RECOMMEND}. Llama's extra review produces a hold without improving restored service.
\sources{Exact status logic: \texttt{src/standard\_agent\_workflow.py}. Verbatim returned reason: \texttt{results/standard\_agent\_workflow.json}, agentic restoration, Explanation call 6.}
\note{Optional. Read the model's actual reason critically. It cites successful checks but concludes REVIEW; the software respects that permitted output. No hidden numerical failure was found in this record, and the fixed comparator used the same numerical tools. The reason is the model's returned explanation, not privileged access to its internal reasoning. A model recommendation in either route still requires human review and is not dispatch.}
\end{frame}

\begin{frame}[t]{Restoration code 12/12: measured time reveals benefit and cost}
\deepdive
\implcode{standard_agent_workflow}{62}{67}
\footnotesize
\begin{tabular}{@{}>{\raggedright\arraybackslash}p{.54\linewidth}rr@{}}
\toprule Same small-MILP protocol & Fixed & Agentic \\
\midrule
Recorded whole-workflow time (s) & 0.009859 & 6.101782 \\
Recorded local Llama calls & 0 & 2 \\
Served / critical served (kW) & 200 / 90 & 200 / 90 \\
Final status & Recommend & Review \\
\bottomrule
\end{tabular}
\par\smallskip\textbf{Two calls:} Planning 2.883081 s; Explanation 3.207768 s. Same solution, extra escalation.
\par\smallskip\textbf{Separate h06 protocol:} 3.160424 s intake + 1.125596 s numerical runtime = 4.286020 s. Llama completed a request handed off by the parser; manual completion time was not measured.
\par\smallskip These are single recorded timings. Compare coverage and incorrect/extra review alongside machine cost.
\sources{Clock recording: \texttt{src/standard\_agent\_workflow.py}. Exact values: \texttt{results/standard\_agent\_workflow.json} and \texttt{results/confirmation.json}, h06.}
\note{Optional. seconds in the excerpt measures the actual local request, parsing and validation span for either success or failure. Whole-workflow time comes from the encompassing timer; individual stage times were not all recorded, so do not invent a per-stage breakdown. The small-MILP pair is a same-protocol descriptive comparison. The h06 intake study is separate and demonstrates task completion with additional machine work. Human labor, utility latency and statistical uncertainty were not measured.}
\end{frame}
\endgroup


\begin{frame}[t]{Restoration: the benefit depends on the stage}
\deepdive
\small
\begin{tabular}{@{}p{.29\linewidth}p{.27\linewidth}p{.34\linewidth}@{}}
\toprule Comparison & Conventional baseline & Agentic workflow \\
\midrule
Six-node MILP, R4 & 200 kW; recommend & Same 200 kW; REVIEW \\
Clear requests, R5 & Correct completion 3/5 & Correct completion 5/5 \\
Missing/prohibited request & Correct clarification 2/2 & Correct clarification 2/2 \\
\bottomrule
\end{tabular}
\takeaway{The observed gain is a more complete language-to-tool handoff. The optimizer supplies the restoration solution; Llama does not improve its optimum and can add an unnecessary review.}
\small\textbf{Operator handoff:} confirm the permitted path, critical demand and shedding authority. Now execute the complete five-stage refinement loop and compare the same-budget search.
\sources{Recorded MILP pair: \texttt{results/standard\_agent\_workflow.json}. Descriptive restoration-language subset: \texttt{results/confirmation.json}. Restoration formulation: \href{https://doi.org/10.1109/TPWRS.2017.2720122}{Chen et al. (2018)}.}
\note{132--135 minutes. Close this use case before introducing cyber response. These are two distinct experiments: a six-node optimization trace and language intake feeding a restricted feeder catalog. Do not combine the 200 kW toy result with the 1.844 MW feeder result. The language subset demonstrates two additional correctly completed requests; it does not establish a superior solver or a field-wide accuracy rate. The six-node trace preserves the actual unnecessary escalation. Both restoration demonstrations use declared trustworthy inputs. The next use case tests what happens when the measurements supporting a calculation are corrupted.}
\end{frame}
